\documentclass[11pt,a4paper]{article}
\usepackage[T1]{fontenc}
\usepackage[utf8]{inputenc}
\usepackage{lmodern}
\DeclareUnicodeCharacter{1E93}{\d{z}}
\usepackage{amsmath,amssymb}
\usepackage[margin=25mm]{geometry}
\usepackage[hidelinks]{hyperref}
\hypersetup{pdftitle={A bounded Alpoge--Sussman peek: what returns to Newton in 0+1},pdfauthor={navstokgap research working notes}}
\newcommand{\tightlist}{\setlength{\itemsep}{0pt}\setlength{\parskip}{0pt}}
\setlength{\emergencystretch}{3em}
\setcounter{secnumdepth}{0}
\begin{document}
\section{A bounded Alpoge--Sussman peek: what returns to Newton in
0+1}\label{a-bounded-alpogesussman-peek-what-returns-to-newton-in-01}

For an exact noncircular Kepler orbit, requiring a local, reversible
Euclidean transfer for its first radial harmonic gives an apparent
action equal to twice the classical radial action. Its second harmonic
is incompatible with that transfer. Keeping just the two orbital speed
bounds already forces a relative discrepancy of at least \(9/80\) in a
real-time squared displacement. A positive jump process reproduces both
frequency laws, and a third radial reading tests positivity across an
inserted cut. All the orbital data come from one Newtonian ellipse.
\hyperref[kepler-motion-and-reflection-at-a-cut]{Kepler motion and
reflection at a cut}.

\textbf{Result and source comparison, 2026-10-27.} The useful immediate
imports are a connected-response test for interaction, retention of the
complete density through blocking, and reflection positivity as a
condition that survives limits. None of these alone selects a positive
\(\hbar\). In our quartic insertion law, an exact moment identity gives
a strictly negative fourth connected response at every positive
\(\hbar\). Its raw magnitude vanishes as a time cell shrinks, while the
correctly scaled quartic response retains the original interaction.

The user authorized this bounded peek after the Newton work and status
checkpoint were pushed. The source is
\href{https://github.com/ethanwsussman/Yang--Mills-LLM-experiment}{ethanwsussman/Yang--Mills-LLM-experiment},
pinned at \texttt{52cee131db140d7467cab4a263a1f04aa15a260f}. The
\href{../docs/AlpogeSussman_YangMills_2026.md}{source companion} records
the downloaded files and exact reading coverage. These are selected
passages, not a proof review of QED3, YM3 or YM4. The repository itself
labels the manuscripts tentative and not satisfactorily vetted.

Retrieval for the calculation: the
\href{newton-nonlinear-time-cell.md}{nonlinear time cell} supplies the
normalized full insertion law and its moment bounds; the
\href{balaban-newton-time-blocking.md}{time-blocking note} separates
finite cell control from a continuum scale; and
\href{refinement-composition-and-limit.md}{refinement composition} keeps
surviving physical information distinct from geometric convergence. The
finite identities below are written derivations, unrefereed, with no
literature priority claim.

\textbf{Targeted check (GPT-6.1 Terra, 2026-10-27).} Terra checked
sections 2--4, including the strict cumulant sign, all \(\hbar\) powers,
the \(T^5\) coefficient and the Gaussian reflection inequality at zero.
Its test-function domain clarification is incorporated below. It did not
review the external manuscripts or physical selection premises.

\subsection{1. What is actually in the current
source}\label{what-is-actually-in-the-current-source}

The current README describes separate QED3 and YM3 manuscripts and a
rotation argument called RP1. It continues to describe YM4 as work being
developed and digested. This is an existence programme; this peek
establishes no mass-gap claim. The following statements are the source's
claims, at the indicated reading levels.

\begin{itemize}
\tightlist
\item
  \textbf{RP1, statement and selected argument passages read.}
  Translation invariance, an exponential moment bound, and reflection
  positivity across coordinate and diagonal mirrors are claimed to imply
  \(O(d)\) covariance of separated-point correlations for \(d\ge2\) and
  tensor rotation weights of absolute value at most two. The conclusion
  allows contact terms at coincident points. Its OS Hilbert-space
  passage constructs a nonnegative time generator; the later rotation
  argument uses a second direction, analytic continuation and angular
  Fourier coefficients. It is not a theorem that positivity in time
  alone produces spacetime rotation symmetry. The root README describes
  RPM1 as a hybrid exposition partly written by Ethan Sussman and
  completed by Astra; it does not specify an Astra version or attribute
  all of the Yang--Mills work to that model.
\item
  \textbf{PR1, abstract and two theorem statements read.} A compact
  probe has a positive fourth logarithmic determinant response. The
  proposed lattice comparison separates a cutoff error, uniform in torus
  size, from an exponentially small volume error. This is a concrete
  nontriviality witness. Its sign belongs to its fermion determinant
  observable; it is not the sign asserted for our scalar moment below.
\item
  \textbf{NL1, abstract, setup and selected theorem clauses read.} The
  nonlinear step keeps local activities, the mass and vacuum terms, and
  an independent external-source scalar term. It claims an exact
  outgoing density and an unnormalized restricted marginal, under
  explicit primitive and domain hypotheses. The methodological lesson is
  to keep the source and scalar terms when differentiating or iterating
  the block map.
\end{itemize}

The dated October 23 revision list was also read. It reports local
corrections, including support conventions, normalization and the
placement of folded cyclic traces. Its assertion that conclusions are
unchanged is not an independent verification here.

\subsection{2. A connected fourth response in the exact Newton
cell}\label{a-connected-fourth-response-in-the-exact-newton-cell}

Use the convex quartic cell at zero source, with

\[
 d\nu(z)=Z^{-1}\exp\left[-\frac1\hbar
                 \left(\frac a2z^2+\frac b4z^4\right)\right]dz,
 \qquad a,b,\hbar>0,
 \qquad m_j=\mathbb E_\nu Z^j.
 \tag{1}
\]

Here \(a=mT/(uv)\) and \(b=w\lambda\), with \(T=u+v\) and \(w=T/2\). The
coefficient \(b\) has units action/length\(^4\). Symmetry gives zero odd
moments, and the quartic tail justifies all differentiations and
integrations by parts below.

Integration by parts against \(z\) and \(z^3\) gives

\[
 a m_2+b m_4=\hbar,\qquad
 a m_4+b m_6=3\hbar m_2.
 \tag{2}
\]

Multiply the second identity by \(m_2\), the first by \(m_4\), and
subtract:

\[
 \boxed{\quad m_4-3m_2^2
       =-\frac b\hbar\left(m_2m_6-m_4^2\right)<0.\quad}
 \tag{3}
\]

Cauchy--Schwarz gives \(m_4^2\le m_2m_6\) by comparing \(Z\) and
\(Z^3\). Equality would require \(Z^3\) proportional to \(Z\) almost
surely, which is impossible for this everywhere positive density. Thus
(3) is a strict finite test of its departure from a Gaussian law. No
small-\(b\) or small-\(\hbar\) expansion was used.

For a source \(j\) of units action/length, let
\(N(j)=\mathbb E e^{jZ/\hbar}\). Then

\[
 (\log N)^{(4)}(0)=\frac{m_4-3m_2^2}{\hbar^4},\qquad
 N^{(4)}(0)=\frac{m_4}{\hbar^4}.
 \tag{4}
\]

The subtraction matters: an ordinary fourth moment is nonzero even for a
free Gaussian. A connected response removes the three Gaussian pairings.
For the retained action of the insertion cell,

\[
 A_\hbar(\mu)-A_\hbar(0)
    =\frac a2\mu^2-\hbar\log N(a\mu),\qquad
 A_\hbar^{(4)}(0)
    =-\frac{a^4}{\hbar^3}(m_4-3m_2^2)>0.
 \tag{5}
\]

This uses the same normalization and source convention as the
nonlinear-cell note. It is an additional full-law response identity, not
an inference from the background trajectory alone.

\subsection{3. Refinement can hide an interaction in an unscaled
probe}\label{refinement-can-hide-an-interaction-in-an-unscaled-probe}

Put \(g=b\hbar/a^2\) and \(Y=\sqrt{a/\hbar}\,Z\). The normalized law of
\(Y\) is proportional to \(e^{-Y^2/2-gY^4/4}\). As \(g\downarrow0\),
dominated convergence gives its second, fourth and sixth moments
\(1,3,15\). Equation (3) therefore gives

\[
 m_4-3m_2^2\sim-\frac{6b\hbar^3}{a^4},\qquad
 A_\hbar^{(4)}(0)\sim6b.
 \tag{6}
\]

For equal time pieces \(u=v=T/2\), \(a=4m/T\), \(b=\lambda T/2\) and
\(g=\lambda\hbar T^3/(32m^2)\). At fixed \(m,\lambda,\hbar>0\),
shrinking the cell consequently gives

\[
 m_4-3m_2^2\sim-\frac{3\lambda\hbar^3}{256m^4}T^5,
 \qquad \frac{A_\hbar^{(4)}(0)}w\longrightarrow6\lambda.
 \tag{7}
\]

Thus the raw local cumulant tends to zero even though the retained
potential's fourth derivative tends to the nonzero derivative of
\(\lambda\mu^4/4\). Neither a finite nonzero cell cumulant nor its raw
vanishing establishes the status of an interacting continuum theory. A
fixed macroscopic observable and an estimate uniform in refinement are
required. That is the useful parallel with PR1's cutoff and volume
comparison, without importing its unreviewed field-theory proof.

\subsection{4. The 0+1 reflection-positivity control, including zero
action}\label{the-01-reflection-positivity-control-including-zero-action}

RP1's rotation theorem needs \(d\ge2\). Its time-positivity ingredient
has a direct \(0+1\) control. Fix \(m,\omega>0\) and consider a centered
stationary Gaussian process with covariance

\[
 C_\hbar(t-s)=\frac{\hbar}{2m\omega}e^{-\omega|t-s|},
 \qquad \hbar\ge0.
 \tag{8}
\]

At zero multiplier this is the deterministic zero process. For real
smooth compactly supported positive-time tests \(f_i\), write
\(Q(f)=\iint f(t)C_\hbar(t-s)f(s)dt\,ds\) and
\(u_i=\int_0^\infty f_i(t)e^{-\omega t}dt\). Time reflection is
\((\theta f)(t)=f(-t)\). Its characteristic functional obeys

\[
 \Phi(f_i-\theta f_j)
 =e^{-Q(f_i)/2}e^{-Q(f_j)/2}
       \exp\left(\frac{\hbar}{2m\omega}u_i u_j\right).
 \tag{9}
\]

For any finite collection, the last kernel is positive semidefinite:
expand its exponential into the sum of rank-one kernels
\([\hbar/(2m\omega)]^n u_i^n u_j^n/n!\). Multiplication by the displayed
positive diagonal factors preserves positivity. This proves the full
Gaussian exponential-observable reflection inequality, including at
\(\hbar=0\), where every matrix entry is one.

Hence reflection positivity and a time generator bounded below do not
exclude a trivial zero-action member. Positive \(\hbar\) in (8) is a
supplied covariance scale; the quartic witness (3) is nonzero for each
positive value but collapses at zero. No universal positive action value
or relativistic \(c\) follows from these tests.

\subsection{Kepler motion and reflection at a
cut}\label{kepler-motion-and-reflection-at-a-cut}

\textbf{6.1 Sol Max, 2026-11-19.} Newton's revised Lemma X uses initial
displacement to determine the force, and Proposition XXXI asks for an
elliptic position at an assigned time. Both use the time of the motion.
To test whether positivity adds a physical premise, keep that time \(t\)
distinct from an auxiliary transfer time \(\tau\).
\href{newton-insertion-action.md\#lemma-x-and-the-measurement-of-initial-force}{Lemma
X, 1687 and 1713} and
\href{../docs/Newton_Principia_Motte1846.md\#position-on-an-ellipse-at-an-assigned-time}{Proposition
XXXI, 1713} (edition passages).

Take a classical bound orbit of mass \(m\) in the potential \(-m\mu/r\),
with \(\mu=GM>0\), energy \(E<0\) and angular momentum \(J>0\). Write
\(A=-m\mu/(2E)\), \(I=m\sqrt{\mu A}\), \(e=\sqrt{1-(J/I)^2}\in(0,1)\)
and \(\nu=\sqrt{\mu/A^3}\). These fix the length and time scales. Under
the stationary preparation the mean anomaly \(\phi\) is uniform on the
circle and advances by \(\nu t\). The radius is \(r=A(1-e\cos u)\),
where \(u-e\sin u=\phi\). On \(0\le u\le\pi\), \(r\) goes from periapsis
to apoapsis. Thus \(X=\cos\phi\) and \(Y=\cos(2\phi)=2X^2-1\) are
bounded functions of radius alone, with the same values on the return
leg. The law of \(X\) is \(\pi_0(dx)=dx/(\pi\sqrt{1-x^2})\) on
\((-1,1)\).
\href{self-sufficiency-contrast.md\#stationary-kepler-preparations}{Stationary
Kepler preparation and source for the orbit clock} (Kepler equations,
passage).

The Newtonian trajectory obeys the exact second-difference identity
\(X_{-t}-2X_0+X_t=2(\cos(\nu t)-1)X_0\). Its normalized limit is
\(-\nu^2X_0\). The actual inward force \(m\mu/r^2\) is finite throughout
the orbit. Yet time reflection \(\Theta:t\mapsto-t\) gives, for
\(F=X_t-X_0\),

\[
 \mathbb E[(\Theta F)F]=\cos(\nu t)\,[\cos(\nu t)-1]<0,
 \qquad 0<\nu t<\pi/2 .                                           \tag{10}
\]

The sign follows from \(\mathbb E[X_0X_t]=\tfrac12\cos(\nu t)\). Moving
the time-zero reading to a positive time tending to zero gives the same
strict sign by continuity. Newtonian time reversal and finite force
therefore admit a negative reflected square at every sufficiently fine
cut.

There is also a bound that uses just one physical time lag. For any
stationary candidate let \(D_f(t)=\mathbb E[(f_t-f_0)^2]\). Positivity
of the reflected square of \(f_t-f_0\) requires \(D_f(2t)\le2D_f(t)\).
With \(T=\pi/(2\nu)\) the ellipse supplies \(D_X(T)=1\). A candidate
preserving this single number and cut positivity at each dyadic step
must have \(D_X(t_n)/t_n\ge1/T\) at \(t_n=T/2^n\). On the orbital
interval \(|dX/dr|\le K\) with \(K=(1+e)^2/(Ae)\), since
\(|\sin(u-e\sin u)-\sin u|\le e\sin u\). Extend \(X\) by its endpoint
values outside that interval. It follows that

\[
 m\,\frac{D_r(t_n)}{t_n}\ \ge\ \frac{m}{K^2T}
 =\frac{2m\nu(Ae)^2}{\pi(1+e)^4}>0 .
\]

The right side has action units. On the actual orbit \(|\dot X|\le\nu\)
and \(D_X(t_n)/t_n\to0\). Thus applying cut positivity in mechanical
time would contradict finite-speed motion already for a single orbit.
Conditionally joining two independent legs at the observed radius would
discard the sign of the radial velocity that Newton carries through the
cut.
\href{conservative-harmonic-receiver.md\#8-what-information-survives-a-cut}{Position
and momentum at a classical cut}.

\subsubsection{What follows from keeping the orbital
frequencies}\label{what-follows-from-keeping-the-orbital-frequencies}

Use a specified additional rule for the auxiliary time. Keep the
positive frequencies in the real cosine correlations and give frequency
\(j\nu\) the transfer factor \(e^{-j\nu\tau}\) for \(\tau\ge0\). This
rule acts on the spectral weights. Substitution of \(i\tau\) for \(t\)
in the real cosine would give a hyperbolic cosine. For
\(f_j=\cos(j\phi)\), the prescribed correlation is
\(C_j(\tau)=\tfrac12e^{-j\nu\tau}\), \(j\ge1\).

There is a classical reflection-positive realization at every
subdivision. On the circle take the Poisson semigroup, whose multiplier
on \(e^{ij\phi}\) is \(b^{|j|}\) with \(b=e^{-\nu\tau}\). Its density
relative to uniform angle is \((1-b^2)/(1-2b\cos\phi+b^2)\). This is a
conservative symmetric Markov transfer. Given the middle phase, its
reflected past and its future are independent with the same conditional
law, so reflected cylinder squares are nonnegative. Its phase makes
jumps. The construction uses the orbital frequency and uniform phase,
with no phase constant.
\href{https://www.math.ru.nl/~mueger/Stein1.pdf}{Stein and Shakarchi,
Chapter 2, §5.4} (Poisson kernel and Abel multipliers, passage).

Now suppose the auxiliary transfer is a conservative symmetric Markov
semigroup \(P_\tau=e^{-\tau H}\) on a stationary probability space, with
a regular, \textbf{strongly local} Dirichlet form \(\mathcal E\). This
is the diffusion hypothesis used here. Additional state variables are
allowed. Suppose \(X\) has the stated radial law and its first mode has
the two correlations \(C_1(s),C_1(2s)\) above at one fixed \(s>0\). Then
\(\|P_sX-e^{-\nu s}X\|^2=0\) and the spectral theorem gives
\(HX=\nu X\). For bounded \(X\) the energy measure \(d\Gamma(X)\) is
normalized by \(\int d\Gamma(X)=\mathcal E(X,X)\). The strongly local
chain rule gives

\[
 \int h'(X)\,d\Gamma(X)
 =\mathcal E(X,h(X))
 =\nu\int xh(x)\,d\pi_0
 =\nu\int(1-x^2)h'(x)\,d\pi_0 .                                  \tag{11}
\]

The last equality is integration by parts with the stated arcsine
density. Taking primitives of continuous tests identifies the
pushforward of the energy measure. Hence for every \(C^1\) function
\(g\) on \([-1,1]\),
\(\mathcal E(g(X),g(X))=\nu\int(1-x^2)g'(x)^2\,d\pi_0\).
\href{https://www.tu-chemnitz.de/mathematik/analysis/preprints/schnol.pdf}{Boutet
de Monvel, Lenz and Stollmann, pp.~192--193} (regular strongly local
forms and energy-measure chain rule, passage).

For \(r=g(X)\) the derivative is finite at both turning points, since
\(dr/d\phi=Ae\sin u/(1-e\cos u)\) vanishes linearly there. Writing
\(I_r=(2\pi)^{-1}\oint p_r\,dr\), (11) gives

\[
 \begin{aligned}
 m\lim_{\tau\downarrow0}\frac{\mathbb E[(r_\tau-r_0)^2]}{\tau}
 &=2m\mathcal E(r,r)
 =2m\nu\int_0^{2\pi}\left(\frac{dr}{d\phi}\right)^2\frac{d\phi}{2\pi}\\
 &=\frac{2m}{\nu}\,\overline{\dot r^{\,2}}=2I_r=2(I-J).
 \end{aligned}                                                     \tag{12}
\]

The dot and overbar in the last line use the mechanical time and its
orbital average. The local circle heat semigroup with generator
\(\nu\partial_\phi^2\) realizes the first mode and (12). At \(I=3J/2\)
the apparent action is \(J\). It vanishes in the circular limit at fixed
\(J\). The local energy density in the physical radius is
\(\nu(dr/d\phi)^2\) times the radial probability. It varies around the
orbit and tends to zero at the turning points. A homogeneous kinetic
coefficient \(a/(2m)\) cannot match it for any constant \(a\).
\href{rivero-1998-conjecture-central-forces.md\#5-hooke-and-kepler-radial-actions-and-maslov-phases}{Kepler
radial action, equation (11)}.

There is an exact test with a second radial function. Equation (11)
gives \(\mathcal E(Y,Y)=16\nu\int(x^2-x^4)d\pi_0=2\nu\). The two
prescribed readings \(C_2(s),C_2(2s)\) instead force \(HY=2\nu Y\),
hence \(\mathcal E(Y,Y)=\nu\), since \(\|Y\|^2=1/2\). Thus four exact
correlations at two fixed positive lags rule out every such strongly
local transfer in the same clock, including hidden variables. The
positive Poisson transfer has the prescribed second mode. The local
circle heat transfer gives \(e^{-4\nu\tau}\) on it.

\subsubsection{A finite physical-time gap from the orbital
speeds}\label{a-finite-physical-time-gap-from-the-orbital-speeds}

The same conflict has a quantitative test in mechanical time if the
transfer is to reconstruct the real, even correlation as
\(\langle f,\cos(tH)f\rangle\). This is a spectral matching premise for
this particular stationary preparation. A Hamiltonian bounded below can
still have transitions to lower energies from an excited state.

On the ellipse \(\overline{\dot X^2}=\nu^2/2\) and
\(\overline{\dot Y^2}=2\nu^2\). For a reconstructed correlation its
negative second derivative at zero is \(\|Hf\|^2\). Impose only the
corresponding upper bounds \(\|HX\|\le\nu/\sqrt2\),
\(\|HY\|\le\sqrt2\,\nu\), with the same equal-time radial law. These are
frequency-moment bounds from physical time, in inverse seconds squared
when squared. They allow rough paths in auxiliary time.

Strong locality alone gives \(\mathcal E(Y,Y)=(16/3)\mathcal E(X,X^3)\)
by the chain rule, without assuming either mode is an eigenfunction. The
second speed bound gives \(\mathcal E(Y,Y)\le\|Y\|\|HY\|\le\nu\). Since
\(\int X^4=3/8\), \(\int X^6=5/16\) in the stationary probability,

\[
 \|(H-\nu)X\|\,\|X^3\|\ge
 \nu\int X^4-\langle HX,X^3\rangle\ge 3\nu/16 .
\]

Squaring and using the first speed bound gives
\(9\nu^2/80\le\|(H-\nu)X\|^2\le\nu^2-2\nu\mathcal E(X,X)\). The mean
frequency in the normalized spectral measure of \(X\) is therefore
\(2\mathcal E(X,X)\le71\nu/80\). In particular
\(2\langle X,P_sX\rangle\ge e^{-71\nu s/80}\) by Jensen's inequality.

For an entirely real-time reading take \(u_*\) maximizing
\((1-\cos u)/u\) on \(u>0\), and put \(t_*=u_*/\nu\). The maximizer is
the root of \(\tan(u_*/2)=u_*\) in \((\pi/2,\pi)\). Integrating the
inequality \(1-\cos z\le z(1-\cos u_*)/u_*\) against the nonnegative
frequency measure gives

\[
 2\langle X,(1-\cos(t_*H))X\rangle
 \le \frac{71}{80}(1-\cos u_*),\qquad
 D_X^{\rm Newton}(t_*)=1-\cos u_* .                               \tag{13}
\]

This discrepancy needs neither exact exponential readings nor a limit of
readout precision. Its hypotheses are stationarity with the specified
radial marginal, the two mechanical speed bounds, strong locality, and
the proposed spectral relation between the clocks.

\subsubsection{What insertion between two readings
tests}\label{what-insertion-between-two-readings-tests}

Smooth paths can reproduce even the \emph{complete} two-time laws of the
Poisson clock. Draw a uniform phase \(\Phi\) and an independent angular
speed \(V\) with Cauchy density \(\nu/[\pi(v^2+\nu^2)]\). Set
\(\phi_\tau=\Phi+V\tau\) modulo \(2\pi\), using the same \(V\)
throughout. The characteristic function
\(\mathbb E e^{ijV\tau}=e^{-|j|\nu|\tau|}\) gives exactly the same pair
transition densities at every lag. This is a random parametrization in
auxiliary time. Its mean squared speed diverges, whereas the physical
Kepler clock has speed \(\nu\).
\href{https://people.math.harvard.edu/~ctm/papers/home/text/class/harvard/154/course/course.pdf}{McMullen,
Probability Theory, Theorem II.6} (Cauchy characteristic function,
passage).

Insert \(Y_0\) between \(X_{-s}\) and \(X_s\). In the smooth process
\(X_{-s}X_s=(\cos 2\Phi+\cos 2Vs)/2\), so the three-time mean is
\(1/4\). In the Poisson process it is \(e^{-2\nu s}/4\), by conditioning
on \(\Phi\) at the cut. In fact \textbf{full reflection positivity
forces the second value} given the Poisson two-time laws. Put
\(\lambda=e^{-\nu s}\) and \(U=X_s-\lambda X_0\). Its reflected square
is zero. Cauchy--Schwarz for the reflected form gives
\(\mathbb E[(\Theta U)Y_0X_s]=0\). The pair law gives
\(\mathbb E[X_0Y_0X_s]=\lambda/4\). Together they fix the three-time
mean. All functions are bounded, and the time-zero variable is also
available as a mean-square limit from positive times. Thus retaining the
same hidden speed on both legs fails positivity on histories despite
matching every pair. This is a test of the composition at a cut in
addition to positivity of the two-point kernel.

The separation persists with finite radius and time resolution. Choose
the two lags \(u,v\) independently in \([s,2s]\) with \(s=(\log2)/\nu\).
For the three marks at \(-u,0,v\), the respective means are
\(e^{-\nu|u-v|}/4\ge1/8\) and \(e^{-\nu(u+v)}/4\le1/16\), by the same
trigonometric identity and conditional expectation. These bounds hold
for any distribution of the lags in the windows. Round each radius in
bins of width \(1/(512K)\) before evaluating \(X,Y\). Since
\(|X|,|Y|\le1\) and \(|2x^2-2z^2|\le4|x-z|\) there, each three-time mean
moves by at most \(6/512\). The resulting intervals remain disjoint. The
equality of all pair laws persists under the identical rounding and
timing.

\subsubsection{Finite accuracy and a local bulk degree of
freedom}\label{finite-accuracy-and-a-local-bulk-degree-of-freedom}

The exact two-mode obstruction has no positive tolerance uniform over
hidden local rates. To see this, add a dimensionless coordinate
\(0\le y\le L\) to the angle. Give the cylinder the probability and form

\[
 d\pi_\eta=\frac{d\phi}{2\pi}\frac{\delta_0(dy)+(\eta/L)\,dy}{1+\eta},
 \qquad
 \mathcal E_\eta(u,u)=\frac{\nu}{1+\eta}
  \int_0^L\int_0^{2\pi}|\nabla u|^2\,\frac{d\phi}{2\pi}\,dy .     \tag{14}
\]

Here \(L,\eta>0\) are dimensionless. Use the periodic \(H^1\) form
domain with its Sobolev trace at \(y=0\). For each fixed \(L,\eta\) the
form norm is equivalent to the \(H^1\) norm by the trace bound. Smooth
functions are dense both in the form domain and uniformly in the
continuous functions on the closed cylinder. This gives a conservative
regular strongly local form, with exactly uniform angle. The
construction uses the harmonic-extension representation of the square
root of a Laplacian.
\href{https://arxiv.org/pdf/math/0608640}{Caffarelli and Silvestre, §1}
(half-space Dirichlet-to-Neumann map, passage).

For \(j\ge1\) set \(h_j=f_j\cosh(j(L-y))/\cosh(jL)\) and
\(\alpha_j=j\nu\tanh(jL)\). The function is harmonic inside and has zero
normal derivative at the top. Integration by parts in (14) gives
\(H_\eta h_j=0\) in the bulk and \(H_\eta h_j=\alpha_j f_j\) at the
bottom. With \(b_\eta=\sqrt{\eta/[2(1+\eta)]}\) this yields
\(\|h_j-f_j\|\le b_\eta\) and
\(\|(H_\eta-\alpha_j)h_j\|\le\alpha_j b_\eta\) in \(L^2(\pi_\eta)\). The
contraction semigroup and Duhamel's formula then give
\(\|P_\tau f_j-e^{-\alpha_j\tau}f_j\|\le3b_\eta\) for all \(\tau\ge0\).
Since \(\|f_j\|=1/\sqrt2\) and
\(|e^{-\alpha_j\tau}-e^{-j\nu\tau}|\le j\nu/\alpha_j-1\),

\[
 \sup_{\tau\ge0,\ j\ge1}
 \left|2\langle f_j,P_\tau f_j\rangle-e^{-j\nu\tau}\right|
 \le 3\sqrt{\frac{\eta}{1+\eta}}+\coth L-1 .                     \tag{15}
\]

Thus local continuous processes approach the target two-time
correlations to arbitrary accuracy as \(\eta\downarrow0\), \(L\uparrow
\infty\). The three-time Poisson mean is approached too, since the
cylinder mean is \(\int Y(P_sX)^2\,d\pi_\eta\) and \(P_sX\) approaches
\(e^{-\nu s}X\) in norm. The bulk has angular and transverse diffusion
coefficient \(\nu L/\eta\). Also \(\mathcal E_\eta(f_j,f_j)=\nu
Lj^2/[2(1+\eta)]\). The diverging initial slopes explain how uniform
agreement of correlations can coexist with a different limit after
division by the time step. In particular this construction cannot keep
the mechanical speed bounds used in (13).

\subsection{5. Consequence for STATE}\label{consequence-for-state}

The authorized Alpoge--Sussman peek is recorded at a fixed commit with
selected source coverage. It contributes a concrete fourth-response test
and a warning to preserve complete source-dependent densities. The
noncircular ellipse supplies a finite force and nonconstant
physical-time motion at every refinement. The apparent action (12)
belongs to an auxiliary diffusion fitted to the first orbital frequency.
Keeping the second frequency calls for a jump transfer or a singular
limit of local transfers. The missing physical premise would have to
relate the clock and its local kinetic coefficient to Newtonian motion
independently of a phase constant. The calculation uses one orbit and
supplies no four-dimensional gauge estimate.
\end{document}
